\chapter{Efectos electrónicos e intermedios reactivos}\label{ch:b1:electronic-effects}

El ácido trifluoroetanoico, \ce{CF3COOH}, es un ácido unas diez mil veces
más fuerte que el ácido etanoico, \ce{CH3COOH}, aunque el enlace O--H
ácido está a tres enlaces de los átomos de flúor. El fenol es un millón de
veces más ácido que el etanol, y una molécula de amoniaco unida a un
anillo de benceno --- la anilina --- es un millón de veces menos básica
que una unida a un grupo metilo. La química orgánica está llena de estos
contrastes, y dos efectos explican la mayoría: la atracción de los
átomos electronegativos a lo largo de los enlaces $\sigma$ y la
deslocalización de los pares de electrones sobre los sistemas $\pi$. Los
mismos dos efectos deciden qué enlaces se rompen, qué intermedios se
forman y dónde atacan los reactivos: son la gramática de los \omterm{def:b1:elementary-steps:elementary-step}{mecanismos
de reacción} de los capítulos siguientes.

\begin{recall}
La \omterm{def:b1:periodicity:electronegativity}{electronegatividad} y sus tendencias (\cref{ch:b1:periodicity}); las
\omterm{def:b1:lewis-resonance-vsepr:lewis}{estructuras de Lewis}, las \omterm{def:b1:lewis-resonance-vsepr:formal-charge}{cargas formales} y la resonancia
(\cref{ch:b1:lewis-resonance-vsepr}); el $\mathrm{p}K_a$ y la fuerza de
los ácidos y las bases (\cref{ch:b1:predominant-reaction}). El volumen
escolar (año 12) dibujó \omterm{def:b1:electronic-effects:curly-arrow}{flechas curvas} de un \omterm{def:b1:electronic-effects:nucleophile}{nucleófilo} a un
\omterm{def:b1:electronic-effects:nucleophile}{electrófilo}.
\end{recall}

\section{Efectos inductivos}

\begin{definition}[Efecto inductivo]\label{def:b1:electronic-effects:inductive}
El \emph{efecto inductivo}\index{efecto inductivo} es el desplazamiento de
la densidad electrónica a lo largo de los enlaces $\sigma$ hacia un átomo
o grupo electronegativo. Un grupo atractor (F, Cl, O, \ce{NO2},
\ce{CF3}) tiene un efecto $-I$; los grupos alquilo y los átomos con carga
negativa empujan electrones, un efecto $+I$. El efecto se debilita
rápidamente con el número de enlaces intermedios.
\end{definition}

\begin{proposition}[Efecto inductivo y acidez]\label{prop:b1:electronic-effects:inductive-pka}
Un grupo atractor cerca del centro ácido de un ácido carboxílico
estabiliza el ion carboxilato al repartir su carga negativa, y rebaja el
$\mathrm{p}K_a$: cuantos más grupos atractores y más cercanos, más fuerte
es el ácido.
\end{proposition}

\begin{proof}
El $\mathrm{p}K_a$ mide la posición de $\mathrm{RCOOH} \rightleftharpoons
\mathrm{RCOO^-} + \ce{H+}$. Un grupo que atrae densidad electrónica del
carboxilato reparte su carga entre más átomos, lo que rebaja su energía
más que la del ácido neutro; el equilibrio se desplaza a la derecha y
$K_a$ aumenta. La serie medida lo confirma.
\end{proof}

\begin{omfigure}
\begin{tikzpicture}
\begin{axis}[
  xbar, width=0.78\linewidth, height=5.2cm, xmin=0, xmax=5.4,
  symbolic y coords={TFA,TCA,DCA,MCA,AA}, ytick=data,
  yticklabels={\ce{CF3COOH}, \ce{CCl3COOH}, \ce{CHCl2COOH}, \ce{CH2ClCOOH}, \ce{CH3COOH}},
  xlabel={$\mathrm{p}K_a$ a \qty{25}{\celsius}}, bar width=9pt,
  nodes near coords, nodes near coords align={horizontal},
  every node near coord/.append style={font=\scriptsize},
  tick label style={font=\small}, label style={font=\small}, enlarge y limits=0.15]
\addplot[fill=omDef!45, draw=omDef] coordinates {(0.52,TFA) (0.51,TCA) (1.35,DCA) (2.87,MCA) (4.76,AA)};
\end{axis}
\end{tikzpicture}

{\small Los ácidos etanoicos halogenados. Cada cloro rebaja el
$\mathrm{p}K_a$; con tres, el ácido es unas diez mil veces más fuerte que
el ácido etanoico. Los ácidos trifluoro- y tricloroetanoico son igual de
fuertes dentro de la incertidumbre de las medidas en ácidos tan
fuertes.}
% ledger: pka:ethanoic, pka:chloroethanoic, pka:dichloroethanoic, pka:trichloroethanoic, pka:trifluoroethanoic
\end{omfigure}

Los grupos alquilo actúan en sentido contrario, un poco: el ácido
propanoico ($\mathrm{p}K_a
= 4.89$) es algo más débil que el ácido etanoico (4.76), y el ácido
metanoico (3.74), sin grupo alquilo, es más fuerte que ambos.
% ledger: pka:propanoic, pka:methanoic

\section{Efectos mesómeros}

Cuando un átomo que lleva un \omterm{def:b1:lewis-resonance-vsepr:lewis}{par no enlazante}, o una carga, está unido a
un sistema $\pi$, sus electrones pueden compartirse con ese sistema; las
\omterm{def:b1:lewis-resonance-vsepr:resonance}{formas resonantes} (\cref{ch:b1:lewis-resonance-vsepr}) describen ese
reparto.

\begin{definition}[Efecto mesómero, grupos dadores y aceptores]\label{def:b1:electronic-effects:mesomeric}
El \emph{efecto mesómero}\index{efecto mesomero@efecto mesómero} es el desplazamiento de la
densidad electrónica a través de un sistema $\pi$ conjugado, que muestran
las \omterm{def:b1:lewis-resonance-vsepr:resonance}{formas resonantes}. Un \emph{grupo dador}\index{grupo dador} cede un
\omterm{def:b1:lewis-resonance-vsepr:lewis}{par no enlazante} al sistema, un efecto $+M$ (\ce{-OH}, \ce{-O^-},
\ce{-NH2}, \ce{-OR}); un \emph{grupo aceptor}\index{grupo aceptor} atrae
electrones de él a través de un enlace $\pi$ con un átomo
electronegativo, un efecto $-M$ (\ce{-NO2},
\ce{-C=O}, \ce{-C#N}).
\end{definition}

\begin{omfigure}
\setchemfig{atom sep=2em}
\schemestart
\chemfig{CH_3-C(=[1]O)-[7]O^{-}}
\arrow{<->}
\chemfig{CH_3-C(-[1]O^{-})=[7]O}
\schemestop
\hspace{2.5em}
\schemestart
\chemfig{*6(=-=(-O^{-})-=-)}
\arrow{<->}
\chemfig{*6((=O)-=-\chemabove{}{\scriptstyle\ominus}-=-)}
\schemestop

\medskip
\schemestart
\chemfig{H_3C-C(=[1]O)-[7]NH_2}
\arrow{<->}
\chemfig{H_3C-C(-[1]O^{-})=[7]NH_2^{+}}
\schemestop

{\small \omterm{def:b1:lewis-resonance-vsepr:resonance}{Formas resonantes} del ion etanoato (la carga, compartida por dos
oxígenos equivalentes), del ion fenóxido (la carga, repartida por el
anillo; se muestra una de las tres estructuras del anillo) y de una amida
(el \omterm{def:b1:lewis-resonance-vsepr:lewis}{par no enlazante} del nitrógeno, compartido con el C=O, lo que hace de
las amidas \omterm{def:b1:predominant-reaction:strong-acid}{bases débiles}).}
\end{omfigure}

\begin{proposition}[Efecto mesómero y acidez]\label{prop:b1:electronic-effects:mesomeric-pka}
Un ácido cuya base conjugada reparte su carga por resonancia es más
fuerte que uno cuya base no puede hacerlo: los ácidos carboxílicos
($\mathrm{p}K_a$ de 4--5) son mucho más fuertes que los alcoholes (en
torno a 16), y el fenol (10.0) está entre ambos. Un \omterm{def:b1:electronic-effects:mesomeric}{grupo aceptor} en
posición \textit{orto} o \textit{para} respecto del OH de un fenol lo
refuerza más que en posición \textit{meta}.
\end{proposition}

\begin{proof}
En un alcóxido, la carga está sobre un solo oxígeno; en un carboxilato,
la comparten dos oxígenos; en un fenóxido, la comparten el oxígeno y tres
carbonos del anillo, menos electronegativos que el oxígeno, una ganancia
menor. Un grupo \ce{NO2} en el carbono \textit{orto} o \textit{para}
puede tomar él mismo la carga mediante una \omterm{def:b1:lewis-resonance-vsepr:resonance}{forma resonante} con
\ce{C=N+(-O^-)-O^-}; en \textit{meta}, ninguna \omterm{def:b1:lewis-resonance-vsepr:resonance}{forma resonante} lleva la
carga al carbono que lleva el \ce{NO2}, y solo queda el \omterm{def:b1:electronic-effects:inductive}{efecto
inductivo}. Los valores medidos: 2-nitrofenol 7.23, 4-nitrofenol 7.16,
3-nitrofenol 8.36, frente a 9.99 para el fenol y 15.9 para el etanol.
\end{proof}
% ledger: pka:ethanol, pka:phenol, pka:2-nitrophenol, pka:3-nitrophenol, pka:4-nitrophenol

El mismo razonamiento se aplica a las bases. La metilamina
($\mathrm{p}K_a$ de su amonio, 10.66) es una base más fuerte que el
amoniaco (9.25, \cref{ch:b1:predominant-reaction}), pues el grupo metilo
empuja electrones hacia el nitrógeno; en la anilina, el \omterm{def:b1:lewis-resonance-vsepr:lewis}{par no enlazante}
se comparte con el anillo y el amonio es mucho más ácido (4.60). La
piridina (5.23) conserva su \omterm{def:b1:lewis-resonance-vsepr:lewis}{par no enlazante}, pero sobre un nitrógeno
retenido en un anillo plano y con más carácter \textit{s}, menos
disponible que en una amina.
% ledger: pka:methylammonium, pka:anilinium, pka:pyridinium, dfg:NH4+_ao

\section{Romper un enlace: los intermedios reactivos}

\begin{definition}[Homólisis y heterólisis]\label{def:b1:electronic-effects:cleavage}
Un \omterm{def:b1:lewis-resonance-vsepr:lewis}{enlace covalente} se rompe por \emph{homólisis}\index{homolisis@homólisis} cuando cada
átomo conserva uno de los dos electrones del enlace, dando dos \omterm{def:b1:electronic-effects:intermediates}{radicales},
y por \emph{heterólisis}\index{heterolisis@heterólisis} cuando un átomo conserva los dos,
dando un catión y un anión.
\end{definition}

\begin{definition}[Intermedios reactivos]\label{def:b1:electronic-effects:intermediates}
Un \emph{intermedio reactivo}\index{intermedio reactivo} es una especie
que se forma en una etapa de un mecanismo y se consume en otra
posterior, demasiado reactiva para acumularse
(\cref{ch:b1:elementary-steps}). Un \emph{carbocatión}\index{carbocation@carbocatión}
tiene un carbono con tres enlaces y una carga positiva (seis \omterm{def:b1:quantum-numbers:configuration}{electrones de
valencia}); un \emph{carbanión}\index{carbanion@carbanión}, un carbono con tres enlaces, un
\omterm{def:b1:lewis-resonance-vsepr:lewis}{par no enlazante} y una carga negativa; un \emph{radical}\index{radical},
un átomo con un \omterm{def:b1:quantum-numbers:unpaired}{electrón desapareado}.
\end{definition}

\begin{definition}[Clase de un carbono]\label{def:b1:electronic-effects:carbon-class}
Un carbono unido a otro carbono es un \emph{carbono
primario}\index{carbono primario}; unido a dos, un \emph{carbono
secundario}\index{carbono secundario}; unido a tres, un \emph{carbono
terciario}\index{carbono terciario}. Un \omterm{def:b1:electronic-effects:intermediates}{carbocatión}, un \omterm{def:b1:electronic-effects:intermediates}{radical} o un
halogenoalcano reciben la clase del carbono correspondiente.
\end{definition}

\begin{omfigure}
\begin{tikzpicture}[font=\small]
  % carbocation: planar, empty p orbital
  \begin{scope}
    \fill[omProp!15, draw=omProp] (0,0.15) ellipse (0.18 and 0.75);
    \fill[omProp!15, draw=omProp] (0,-0.15) ellipse (0.18 and 0.75);
    \draw[thick] (0,0) -- (-1.0,-0.25) node[left] {R};
    \draw[thick] (0,0) -- (0.95,-0.35) node[right] {R};
    \draw[thick] (0,0) -- (0.35,0.45) node[above right] {R};
    \node[fill=white, inner sep=1pt] at (0,0) {\ce{C+}};
    \node at (0,-1.75) {carbocatión: plano,};
    \node at (0,-2.15) {orbital \textit{p} vacío};
  \end{scope}
  % radical: nearly planar, one electron
  \begin{scope}[xshift=4.6cm]
    \draw[omMeth] (0,0.15) ellipse (0.18 and 0.75);
    \draw[omMeth] (0,-0.15) ellipse (0.18 and 0.75);
    \fill (0,0.55) circle (0.05);
    \draw[thick] (0,0) -- (-1.0,-0.25) node[left] {R};
    \draw[thick] (0,0) -- (0.95,-0.35) node[right] {R};
    \draw[thick] (0,0) -- (0.35,0.45) node[above right] {R};
    \node[fill=white, inner sep=1pt] at (0,0) {C};
    \node at (0,-1.75) {radical: casi plano,};
    \node at (0,-2.15) {un electrón};
  \end{scope}
  % carbanion: pyramidal with a lone pair
  \begin{scope}[xshift=9.2cm]
    \draw[thick] (0,0) -- (-0.95,-0.55) node[left] {R};
    \draw[thick] (0,0) -- (0.95,-0.55) node[right] {R};
    \draw[thick] (0,0) -- (0.2,-0.9) node[below] {R};
    \fill[omDef!15, draw=omDef] (0,0.55) ellipse (0.22 and 0.45);
    \fill (-0.07,0.65) circle (0.045); \fill (0.07,0.65) circle (0.045);
    \node[fill=white, inner sep=1pt] at (0,0) {\ce{C-}};
    \node at (0,-1.75) {carbanión: piramidal,};
    \node at (0,-2.15) {par no enlazante};
  \end{scope}
\end{tikzpicture}

{\small Geometrías de los tres \omterm{def:b1:electronic-effects:intermediates}{intermedios reactivos} del carbono (R: H o
un grupo alquilo). El \omterm{def:b1:electronic-effects:intermediates}{carbocatión} plano puede ser atacado por cualquiera
de sus caras, algo que importa para la estereoquímica del
\cref{ch:b1:nucleophilic-substitution}.}
\end{omfigure}

\begin{proposition}[Estabilidad de los carbocationes]\label{prop:b1:electronic-effects:carbocation-order}
Los \omterm{def:b1:electronic-effects:intermediates}{carbocationes} se estabilizan con los grupos que les ceden electrones:
terciario $>$ secundario $>$ primario $>$ metilo; un \omterm{def:b1:electronic-effects:intermediates}{carbocatión} contiguo
a un enlace doble (alílico) o a un anillo de benceno (bencílico) se
estabiliza además por resonancia, y uno contiguo a un átomo con un \omterm{def:b1:lewis-resonance-vsepr:lewis}{par no
enlazante} (\ce{R-O-CH2+}), aún más.
\end{proposition}

\begin{proof}
Cada grupo alquilo empuja densidad electrónica hacia el carbono pobre en
electrones (efecto $+I$, y una interacción relacionada de sus enlaces
C--H con el orbital vacío, descrita en el volumen del segundo año). Un
catión alílico tiene dos \omterm{def:b1:lewis-resonance-vsepr:resonance}{formas resonantes} que sitúan la carga en uno u
otro carbono extremo; un catión bencílico tiene cuatro, tres en el
anillo; un \omterm{def:b1:lewis-resonance-vsepr:lewis}{par no enlazante} del oxígeno forma un cuarto enlace con el
carbono catiónico, \ce{R-O+=CH2}, en el que todos los átomos tienen
octeto.
\end{proof}

Los \omterm{def:b1:electronic-effects:intermediates}{carbaniones} siguen el orden contrario (metilo $>$ primario $>$
secundario $>$ terciario) y se estabilizan con \omterm{def:b1:electronic-effects:mesomeric}{grupos aceptores}; los
\omterm{def:b1:electronic-effects:intermediates}{radicales} siguen el mismo orden que los \omterm{def:b1:electronic-effects:intermediates}{carbocationes}, de forma menos
marcada.

\section{Flechas curvas, nucleófilos y electrófilos}

\begin{definition}[Flechas curvas]\label{def:b1:electronic-effects:curly-arrow}
En un mecanismo, una \emph{flecha curva}\index{flecha curva} muestra el
movimiento de un par de electrones: parte de un \omterm{def:b1:lewis-resonance-vsepr:lewis}{par no enlazante} o de un
enlace y termina en un átomo (donde se forma un \omterm{def:b1:lewis-resonance-vsepr:lewis}{par no enlazante} o un
enlace nuevo) o entre dos átomos (un enlace nuevo). Una \emph{semiflecha
curva}\index{semiflecha curva} muestra el movimiento de un solo electrón,
en las reacciones radicalarias.
\end{definition}

\begin{method}[Dibujar flechas curvas]\label{met:b1:electronic-effects:curly-arrows}
\begin{enumerate}
  \item Se parte de electrones: un \omterm{def:b1:lewis-resonance-vsepr:lewis}{par no enlazante}, un enlace $\pi$ o un
        enlace $\sigma$, nunca de una carga positiva ni de un átomo sin
        electrones.
  \item Se termina donde va el par: un enlace nuevo entre los dos átomos,
        o un \omterm{def:b1:lewis-resonance-vsepr:lewis}{par no enlazante} sobre el átomo más electronegativo cuando
        se rompe un enlace.
  \item Se cuentan las cargas tras cada etapa: el átomo que cede un \omterm{def:b1:lewis-resonance-vsepr:lewis}{par no
        enlazante} a un enlace nuevo gana $+1$, y el que recibe un par
        como \omterm{def:b1:lewis-resonance-vsepr:lewis}{par no enlazante} gana $-1$; la carga total se conserva.
  \item Nunca se supera el octeto en C, N, O: cuando llega un par, otro
        debe irse.
\end{enumerate}
\end{method}

\begin{omfigure}
\setchemfig{atom sep=2.2em}
\schemestart
\chemfig{H_3C-C(-[2]H)(-[6]H)-[@{b}]@{br}Br}
\arrow{->[heterólisis]}[,1.5]
\chemfig{H_3C-C(-[2]H)(-[6]H)^{+}}\+\chemfig{Br^{-}}
\schemestop
\chemmove{\draw[->, shorten <=2pt, shorten >=2pt](b) .. controls +(90:6mm) and +(90:6mm) .. (br);}

\medskip
\schemestart
\chemfig{H@{o}O^{-}}\+\chemfig{@{h}H-@{oc}O-H}
\arrow{<=>}
\chemfig{H_2O}\+\chemfig{HO^{-}}
\schemestop
\chemmove{\draw[->, thick, shorten <=2pt, shorten >=2pt](o) .. controls +(-80:8mm) and +(-100:8mm) .. (h);}
\par\vspace{5mm}

{\small \omterm{def:b1:electronic-effects:curly-arrow}{Flechas curvas}. Arriba, \omterm{def:b1:electronic-effects:cleavage}{heterólisis} del enlace C--Br: el par va al
bromo. Abajo, una transferencia de protón del agua al hidróxido, escrita
como el ataque de un \omterm{def:b1:lewis-resonance-vsepr:lewis}{par no enlazante} de la base sobre el hidrógeno (el
enlace O--H del agua se rompe entonces y su par queda en el oxígeno; la
segunda flecha se deja para el \cref{exo:b1:electronic-effects:4}).}
\end{omfigure}

\begin{definition}[Nucleófilo, electrófilo, grupo saliente]\label{def:b1:electronic-effects:nucleophile}
Un \emph{nucleófilo}\index{nucleofilo@nucleófilo} es una especie que cede un par de
electrones a un átomo distinto del hidrógeno para formar un enlace (una
\omterm{def:b1:lewis-resonance-vsepr:lewis-acid}{base de Lewis}); un \emph{electrófilo}\index{electrofilo@electrófilo} es una especie que
acepta ese par (un \omterm{def:b1:lewis-resonance-vsepr:lewis-acid}{ácido de Lewis}, a menudo un carbono que lleva una
carga parcial positiva). La \emph{nucleofilia}\index{nucleofilia} de una
especie es su reactividad como nucleófilo, medida por \omterm{def:b1:rate-laws:order}{constantes de
velocidad}. Un \emph{grupo saliente}\index{grupo saliente} es el grupo que
se va con el par de enlace cuando un enlace con el carbono se rompe de
forma heterolítica.
\end{definition}

\begin{proposition}[Nucleofilia y basicidad]\label{prop:b1:electronic-effects:nucleophilicity-basicity}
Para \omterm{def:b1:electronic-effects:nucleophile}{nucleófilos} que atacan a través del mismo átomo, la \omterm{def:b1:electronic-effects:nucleophile}{nucleofilia} sigue
a la basicidad: \ce{CH3O-} $>$ \ce{OH-} $>$ \ce{CH3COO-} $>$ \ce{H2O}. Deja
de cumplirse a lo largo de una columna de la tabla periódica en los
\omterm{def:b1:intermolecular-forces-solvents:solvent-classes}{disolventes próticos}, donde los \omterm{def:b1:electronic-effects:nucleophile}{nucleófilos} grandes, polarizables y poco
solvatados reaccionan más deprisa aunque sean bases más débiles
(\ce{I-} $>$ \ce{Br-} $>$ \ce{Cl-} $>$ \ce{F-}), y para las bases
voluminosas, que son malos \omterm{def:b1:electronic-effects:nucleophile}{nucleófilos}.
\end{proposition}

\begin{proof}
La basicidad mide el equilibrio del ataque a un protón; la \omterm{def:b1:electronic-effects:nucleophile}{nucleofilia},
la velocidad del ataque a un carbono; ambas crecen con la disponibilidad
del par de electrones, de ahí el paralelismo para un mismo átomo. La
velocidad depende también de la energía necesaria para despojar al
\omterm{def:b1:electronic-effects:nucleophile}{nucleófilo} del disolvente y de la deformación de su nube electrónica
hacia un carbono impedido: los aniones pequeños como \ce{F-} están
fuertemente unidos a las moléculas de un \omterm{def:b1:intermolecular-forces-solvents:solvent-classes}{disolvente prótico}
(\cref{ch:b1:intermolecular-forces-solvents}), y los grupos voluminosos
dificultan el acercamiento al carbono pero no a un protón.
\end{proof}

Un buen \omterm{def:b1:electronic-effects:nucleophile}{grupo saliente} es una \omterm{def:b1:predominant-reaction:strong-acid}{base débil}, estable una vez que lleva el
par: \ce{I-}, \ce{Br-}, \ce{Cl-}, el agua, los sulfonatos; \ce{OH-} y
\ce{RO-}, \omterm{def:b1:predominant-reaction:strong-acid}{bases fuertes}, son malos \omterm{def:b1:electronic-effects:nucleophile}{grupos salientes} --- la razón por la
que los alcoholes deben activarse antes de una sustitución
(\cref{ch:b1:alcohol-activation}).

\section{Ejercicios}

\begin{exercise}[$\star$]\label{exo:b1:electronic-effects:1}
Clasifíquese cada carbono del 2-metilbutano y del
2-bromo-2-metilpropano como primario, secundario, terciario o cuaternario
(unido a cuatro carbonos). ¿Qué \omterm{def:b1:electronic-effects:intermediates}{carbocatión} se forma con más facilidad a
partir de cada bromuro de fórmula \ce{C4H9Br}?
\end{exercise}

\begin{exercise}[$\star$]\label{exo:b1:electronic-effects:2}
A partir de los valores de $\mathrm{p}K_a$, calcúlese la razón de las
\omterm{def:b1:predominant-reaction:acidity-constant}{constantes de acidez} de los ácidos cloroetanoico y etanoico, y de los
ácidos trifluoroetanoico y etanoico.
\end{exercise}

\begin{exercise}[$\star$]\label{exo:b1:electronic-effects:3}
Dibújense las \omterm{def:b1:lewis-resonance-vsepr:resonance}{formas resonantes} del ion etanoato, del ion 4-nitrofenóxido
(con la carga en un oxígeno del grupo nitro) y del catión alilo
\ce{CH2=CH-CH2+}.
\end{exercise}

\begin{exercise}[$\star$]\label{exo:b1:electronic-effects:4}
Complétense las \omterm{def:b1:electronic-effects:curly-arrow}{flechas curvas} de la transferencia de protón de la
figura y de \ce{NH3 + H3O+ -> NH4+ + H2O}. Compruébense las cargas.
\end{exercise}

\begin{exercise}[$\star\star$]\label{exo:b1:electronic-effects:5}
Ordénense por acidez creciente: el etanol, el fenol, el ácido etanoico,
el 4-nitrofenol, el ácido tricloroetanoico y el agua ($\mathrm{p}K_a$ del
par \ce{H2O}/\ce{OH-}, 14.0). Justifíquese cada paso con un argumento
inductivo o mesómero.
\end{exercise}

\begin{exercise}[$\star\star$]\label{exo:b1:electronic-effects:6}
Ordénense por basicidad creciente: la anilina, la metilamina, el
amoniaco y la piridina. Explíquese la posición de la anilina con una
\omterm{def:b1:lewis-resonance-vsepr:resonance}{forma resonante}.
\end{exercise}

\begin{exercise}[$\star\star$]\label{exo:b1:electronic-effects:7}
¿Por qué el 3-nitrofenol es más débil que el 4-nitrofenol pero sigue
siendo más fuerte que el fenol? Dibújense las \omterm{def:b1:lewis-resonance-vsepr:resonance}{formas resonantes} que
justifican ambas afirmaciones.
\end{exercise}

\begin{exercise}[$\star\star$]\label{exo:b1:electronic-effects:8}
Ordénense los \omterm{def:b1:electronic-effects:intermediates}{carbocationes} \ce{CH3+}, \ce{(CH3)3C+}, \ce{CH2=CH-CH2+},
\ce{CH3CH2+}, \ce{C6H5CH2+} y \ce{CH3OCH2+}, y justifíquese.
\end{exercise}

\begin{exercise}[$\star\star$]\label{exo:b1:electronic-effects:9}
Identifíquense el \omterm{def:b1:electronic-effects:nucleophile}{nucleófilo} y el \omterm{def:b1:electronic-effects:nucleophile}{electrófilo} en: \ce{CH3Br + OH-};
\ce{H2O + H+}; el etanal con un ion cianuro; \ce{BF3 + NH3}. Dibújense las
\omterm{def:b1:electronic-effects:curly-arrow}{flechas curvas}.
\end{exercise}

\begin{exercise}[$\star\star\star$]\label{exo:b1:electronic-effects:10}
Una disolución contiene \qty{0.010}{mol/L} de ácido tricloroetanoico. ¿Es
válida la fórmula del \omterm{def:b1:predominant-reaction:strong-acid}{ácido débil},
$\mathrm{pH} = \frac12(\mathrm{p}K_a - \log C)$? Calcúlense correctamente
el pH y el \omterm{def:b1:predominant-reaction:dissociation}{grado de disociación}.
\end{exercise}

\begin{exercise}[$\star\star\star$]\label{exo:b1:electronic-effects:11}
Los $\mathrm{p}K_a$ de los alcoholes medidos en agua aumentan del metanol
(15.5) al etanol (15.9), al propan-2-ol (17.1) y al 2-metilpropan-2-ol
(19.2). Dense dos explicaciones, una mediante el \omterm{def:b1:electronic-effects:inductive}{efecto inductivo} de los
grupos alquilo sobre el alcóxido y otra mediante la \omterm{def:b1:intermolecular-forces-solvents:dissolution}{solvatación} del
alcóxido.
% ledger: pka:methanol, pka:ethanol, pka:2-propanol, pka:tBuOH
\end{exercise}

\begin{exercise}[$\star\star\star$]\label{exo:b1:electronic-effects:12}
Explíquese por qué las amidas no son ni básicas en el nitrógeno ni
nucleófilas, mientras que las aminas son ambas cosas, con la \omterm{def:b1:lewis-resonance-vsepr:resonance}{forma
resonante} de la amida. ¿Qué átomo de una amida se protona en \omterm{def:b1:predominant-reaction:strong-acid}{ácido
fuerte}?
\end{exercise}

\section{Problema: por qué el ácido trifluoroetanoico es tan fuerte}

\begin{problem}[{Problema de fin de semana --- la serie inductiva de los
ácidos halogenados, carboxilatos frente a alcóxidos, los tres
nitrofenoles, un cálculo por la reacción predominante y la razón de las
constantes de acidez de los ácidos trifluoroetanoico y etanoico}]
\label{pb:b1:electronic-effects:1}
Datos ($\mathrm{p}K_a$ a \qty{25}{\celsius}): ácido etanoico 4.76,
cloroetanoico 2.87, dicloroetanoico 1.35, tricloroetanoico 0.51,
trifluoroetanoico 0.52, etanol 15.9, fenol 9.99, 2-nitrofenol 7.23,
3-nitrofenol 8.36, 4-nitrofenol 7.16; $\mathrm{p}K_e = 14.00$.

\medskip
\noindent\textbf{Parte I --- La serie inductiva.}
\begin{enumerate}
  \item Calcúlese la razón de las $K_a$ en cada paso del ácido etanoico al
        tricloroetanoico.
  \item ¿Es igual el efecto de cada cloro? Coméntese.
  \item Explíquese la tendencia con el \omterm{def:b1:electronic-effects:inductive}{efecto inductivo}.
  \item ¿Por qué un cloro en el carbono de una cadena más larga más
        alejado del COOH casi no tendría efecto?
  \item El flúor es más electronegativo que el cloro. ¿Qué muestra la
        comparación de los ácidos trifluoro- y tricloro-, y por qué es
        difícil medir con precisión esos $\mathrm{p}K_a$?
  \item ¿En qué forma está cada uno de estos ácidos a pH 3?
\end{enumerate}

\medskip
\noindent\textbf{Parte II --- Carboxilatos y alcóxidos.}
\begin{enumerate}[resume]
  \item Dibújense las dos \omterm{def:b1:lewis-resonance-vsepr:resonance}{formas resonantes} del ion etanoato.
  \item ¿Qué predicen para las longitudes de los dos enlaces C--O del ion?
  \item ¿Por qué el ion etóxido no tiene estructuras equivalentes?
  \item Calcúlese la razón de las $K_a$ del ácido etanoico y del etanol.
  \item Sitúese el fenol entre ambos y justifíquese con las \omterm{def:b1:lewis-resonance-vsepr:resonance}{formas
        resonantes} del fenóxido.
  \item ¿Cuáles de entre el etanol, el fenol y el ácido etanoico reaccionan
        casi totalmente con una disolución de hidrogenocarbonato
        ($\mathrm{p}K_a$ 6.37 para
        \ce{CO2,H2O}/\ce{HCO3-})?
\end{enumerate}

\medskip
\noindent\textbf{Parte III --- Los nitrofenoles.}
\begin{enumerate}[resume]
  \item Ordénense por acidez los tres nitrofenoles y el fenol.
  \item Dibújese una \omterm{def:b1:lewis-resonance-vsepr:resonance}{forma resonante} del 4-nitrofenóxido con la carga en
        el grupo nitro.
  \item Demuéstrese que no existe una estructura así para el
        3-nitrofenóxido.
  \item ¿Por qué el 3-nitrofenol sigue siendo más ácido que el fenol?
  \item El 2-nitrofenol es algo más débil que el 4-nitrofenol aunque ambos
        tienen el \omterm{def:b1:electronic-effects:mesomeric}{efecto mesómero}; propóngase una explicación en la que
        intervengan su OH y el grupo nitro cercano.
  \item A pH 7.5, ¿qué nitrofenoles están mayoritariamente en forma de
        anión? ¿Por qué hace eso del 4-nitrofenol un \omterm{def:b1:titration-methods:indicator}{indicador coloreado}?
\end{enumerate}

\medskip
\noindent\textbf{Parte IV --- El ácido trifluoroetanoico en el agua.}
\begin{enumerate}[resume]
  \item Calcúlese el pH del ácido trifluoroetanoico de \qty{0.010}{mol/L}
        (sin suponer una disociación débil).
  \item Calcúlese el pH del ácido etanoico de \qty{0.010}{mol/L}.
  \item Escríbase la reacción del ácido trifluoroetanoico con el etanoato
        de sodio y calcúlese su constante.
  \item Dese la razón $K_a(\ce{CF3COOH})/K_a(\ce{CH3COOH})$, como potencia
        de diez y como número.
\end{enumerate}
\end{problem}
